% =========================================================
% frontmatter/learning-outcomes.tex
% Capaian pembelajaran buku: CPL, CPMK, Sub-CPMK, dan
% pemetaannya ke bab.
% =========================================================

\chapter*{Capaian Pembelajaran Buku}
\addcontentsline{toc}{chapter}{Capaian Pembelajaran Buku}
\markboth{Capaian Pembelajaran Buku}{Capaian Pembelajaran Buku}

Buku ajar ini disusun berdasarkan Rencana Pembelajaran Semester (RPS)
berbasis \textit{Outcome-Based Education}. Seluruh bab diarahkan pada
pencapaian Capaian Pembelajaran Lulusan (CPL) yang dibebankan pada
mata kuliah ini.

\section*{1. Capaian Pembelajaran Lulusan (CPL)}

\begin{itemize}
    \item \textbf{CPL02:} Memiliki pengetahuan yang memadai terkait
          dengan cara kerja sistem komputer dan mampu merancang dan
          mengembangkan berbagai algoritma metode untuk memecahkan
          masalah.

    \item \textbf{CPL03:} Kemampuan merancang dan menganalisis
          algoritma untuk menyelesaikan permasalahan organisasi
          secara optimal, serta memilih dan menerapkannya pada
          bahasa pemrograman tertentu.
\end{itemize}

\section*{2. Capaian Pembelajaran Mata Kuliah (CPMK)}

\begin{longtable}{@{}p{1.7cm}p{9.9cm}p{2.8cm}@{}}
\toprule
\textbf{Kode} & \textbf{Rumusan CPMK} & \textbf{Mendukung} \\
\midrule
\endfirsthead
\toprule
\textbf{Kode} & \textbf{Rumusan CPMK} & \textbf{Mendukung} \\
\midrule
\endhead
\bottomrule
\endfoot
\CPMK{01} & Menjelaskan arsitektur register dan manajemen memori
level rendah. & CPL02, CPMK01 \\
\addlinespace
\CPMK{02} & Memahami hubungan instruksi mesin dengan \textit{hardware}.
& CPL02, CPMK02 \\
\addlinespace
\CPMK{03} & Menulis kode assembly untuk operasi aritmatika dan logika.
& CPL03, CPMK03 \\
\addlinespace
\CPMK{04} & Melakukan \textit{debugging} pada level instruksi memori.
& CPL03, CPMK04 \\
\bottomrule
\end{longtable}

\section*{3. Sub-CPMK dan Pemetaannya ke Bab}

Tabel berikut menunjukkan Sub-CPMK yang harus dicapai mahasiswa,
indikator ketercapaiannya, serta bab tempat kompetensi tersebut
dibangun dan diukur.

\begingroup
\small
\begin{longtable}{@{}p{1.25cm}p{6.9cm}p{1.6cm}p{3.4cm}@{}}
\toprule
\textbf{Sub-CPMK} & \textbf{Indikator Pencapaian} &
\textbf{CPMK} & \textbf{Bab} \\
\midrule
\endfirsthead
\toprule
\textbf{Sub-CPMK} & \textbf{Indikator Pencapaian} &
\textbf{CPMK} & \textbf{Bab} \\
\midrule
\endhead
\bottomrule
\endfoot
1.1 & Menjelaskan arsitektur dasar mikroprosesor 8086
(\textit{Execution Unit} dan \textit{Bus Interface Unit}), peran bus
alamat/data/kontrol, serta tahapan siklus instruksi
(\textit{fetch--decode--execute}). & CPMK01 &
Bab~III, Bab~IV \\
\addlinespace
1.2 & Mengidentifikasi fungsi register umum (AX, BX, CX, DX),
register segmen (CS, DS, SS, ES), register \textit{pointer}/\textit{index}
(SP, BP, SI, DI), dan register \textit{flag} beserta pengaruhnya
terhadap eksekusi program. & CPMK01 & Bab~V \\
\addlinespace
1.3 & Menerapkan konsep segmentasi memori: menghitung alamat fisik
20-bit dari pasangan \textit{segment:offset}, menjelaskan paragraf
16-byte, dan menganalisis skema pengalamatan 1 MB. & CPMK01 & Bab~VI \\
\addlinespace
2.1 & Menjelaskan format instruksi mesin (\textit{opcode}, byte
ModR/M, \textit{displacement}/\textit{immediate}) dan mengaitkannya
dengan operasi yang dieksekusi perangkat keras. & CPMK02 & Bab~VII \\
\addlinespace
2.2 & Menganalisis mode pengalamatan (\textit{immediate},
\textit{direct}, \textit{register indirect}, \textit{indexed},
\textit{base-relative}) dan menentukan segmen \textit{default} yang
digunakan. & CPMK02 & Bab~VIII \\
\addlinespace
2.3 & Menggunakan direktif assembler (DB, DW, DD, EQU, ORG, SEGMENT)
dan menjelaskan alur translasi \textit{source} $\rightarrow$
\textit{object} $\rightarrow$ \textit{executable} oleh assembler dan
linker. & CPMK02 & Bab~IX \\
\addlinespace
2.4 & Menjelaskan mekanisme interupsi dan layanan BIOS/DOS (INT 10h,
INT 21h) sebagai jembatan antara instruksi mesin dan perangkat
keras/I/O. & CPMK02 & Bab~XIV \\
\addlinespace
3.1 & Menulis program assembly lengkap menggunakan instruksi
transfer data (MOV, XCHG, LEA, PUSH/POP) dan instruksi aritmatika
(ADD, SUB, INC, DEC, MUL, IMUL, DIV, IDIV). & CPMK03 & Bab~X \\
\addlinespace
3.2 & Menerapkan operasi logika \textit{bitwise} dan operasi
pergeseran/rotasi (AND, OR, XOR, NOT, TEST, SHL, SHR, SAL, SAR, ROL,
ROR) untuk manipulasi bit. & CPMK03 & Bab~XI \\
\addlinespace
3.3 & Mengimplementasikan struktur kontrol (CMP, TEST, lompatan
bersyarat Jcc, LOOP, dan prosedur dengan CALL/RET) dalam program
assembly. & CPMK03 & Bab~XII \\
\addlinespace
3.4 & Mengembangkan prosedur/subrutin dengan mekanisme
\textit{parameter passing} dan mengelola \textit{stack frame} (PUSH,
POP, BP, konvensi pemanggilan). & CPMK03 & Bab~XIII \\
\addlinespace
4.1 & Melakukan penelusuran program pada level instruksi
(\textit{single-step}, \textit{trace}, \textit{breakpoint}) dan
memantau perubahan register serta \textit{flag}. & CPMK04 & Bab~XV \\
\addlinespace
4.2 & Memeriksa dan menganalisis isi memori dan \textit{stack}
(\textit{dump}) untuk menemukan kesalahan logika, kesalahan
segmentasi, dan \textit{overflow}. & CPMK04 & Bab~XV \\
\addlinespace
4.3 & Menganalisis kode mesin hasil \textit{disassembly} dan
menerapkan strategi optimasi program pada level instruksi. & CPMK04 &
Bab~XVI \\
\bottomrule
\end{longtable}
\endgroup

\section*{4. Bobot Penilaian}

\begin{longtable}{@{}p{3.4cm}p{7.4cm}p{2.3cm}c@{}}
\toprule
\textbf{Komponen} & \textbf{Indikator/CPMK} & \textbf{Sub-CPMK} &
\textbf{Bobot} \\
\midrule
\endfirsthead
\toprule
\textbf{Komponen} & \textbf{Indikator/CPMK} & \textbf{Sub-CPMK} &
\textbf{Bobot} \\
\midrule
\endhead
\bottomrule
\endfoot
Tugas Individu & CPMK01, CPMK02 & 1.1--1.3, 2.1--2.3 & 15\% \\
\addlinespace
Praktikum & CPMK03, CPMK04 & 3.1--3.4, 4.1 & 20\% \\
\addlinespace
Kuis & CPMK01, CPMK02 & 1.2, 1.3, 2.1 & 10\% \\
\addlinespace
UTS (Minggu 8) & CPMK01, CPMK02 & 1.1--2.3 & 20\% \\
\addlinespace
Proyek Akhir & CPMK03, CPMK04 & 3.1--4.3 & 15\% \\
\addlinespace
UAS (Minggu 16) & CPMK01--CPMK04 & 1.1--4.3 & 20\% \\
\midrule
\textbf{Total} & & & \textbf{100\%} \\
\bottomrule
\end{longtable}

\vspace{0.4cm}

\begin{infobox}
Penekanan penilaian mata kuliah ini adalah \textbf{70\% praktik} dan
\textbf{30\% teori}. Komponen praktik menilai kemampuan menulis,
merakit, menjalankan, dan melakukan \textit{debugging} program
assembly; komponen teori menilai pemahaman konsep arsitektur,
instruksi mesin, dan organisasi memori.
\end{infobox}

\clearpage
